\subsection{Fibras formales de los anillos noetherianos henselianos.}
\label{subsection:IV.18.9}

\begin{theorem}[18.9.1]
\label{IV.18.9.1}
Sea $A$ un anillo local noetheriano henseliano cuyas fibras formales son geométricamente normales. Entonces las fibras formales de $A$ son geométricamente íntegras (y, por tanto, geométricamente conexas).
\end{theorem}

Las dos afirmaciones del enunciado son, de hecho, equivalentes, pues todo preesquema localmente noetheriano, conexo y normal es íntegro. Como toda $A$-álgebra íntegra finita es un anillo local henseliano (18.5.9 y 18.5.10), se deduce de (7.3.16.2) que la cuestión se reduce a demostrar que, si además se supone que $A$ es \emph{íntegro}, entonces la fibra del morfismo $\Spec(\widehat A)\to\Spec(A)$ en el punto genérico de $\Spec(A)$ es \emph{íntegra}, lo que se deduce del

\begin{corollary}[18.9.2]
\label{IV.18.9.2}
Sea $A$ un anillo local noetheriano henseliano íntegro cuyas fibras formales son geométricamente normales. Entonces $\widehat A$ es íntegro.
\end{corollary}

En efecto, se deduce de (18.8.16) que $A$ es unirrama, por lo que el teorema se deduce de la hipótesis formulada sobre las fibras formales de $A$ y de (7.6.3).

\begin{corollary}[18.9.3]
\label{IV.18.9.3}
Bajo las hipótesis de (18.9.2), el cuerpo de fracciones $L$ de $\widehat A$ es una extensión separable del cuerpo de fracciones $K$ de $A$, y $K$ es algebraicamente cerrado en $L$.
\end{corollary}

Esto se deduce de que la fibra del morfismo $\Spec(\widehat A)\to\Spec(A)$ en el punto genérico de $\Spec(A)$ es geométricamente íntegra, y de (4.3.2) y (4.3.5).

\begin{corollary}[18.9.4]
\label{IV.18.9.4}
Sea $A$ un anillo local noetheriano y henseliano cuyas fibras formales sean geométricamente normales (por ejemplo, un anillo local noetheriano henseliano y excelente), y pongamos $A'=\widehat A$. Sean $X$ un $A$-preesquema, $X'=X\otimes_A A'$, y $g:X'\to X$ la proyección canónica. Entonces la aplicación $U\mapsto g^{-1}(U)$ es una biyección del conjunto de las partes a la vez abiertas y cerradas de $X$ sobre el conjunto de las partes a la vez abiertas y cerradas de $X'$.
\end{corollary}

Como $g$ es un morfismo fielmente plano y cuasicompacto, se sabe por (2.3.12) que la topología de $X$ es \emph{cociente} de la de $X'$ por la relación de equivalencia definida por $g$. Todo se reduce, por tanto, a ver que una parte abierta y cerrada $U'$ de $X'$ está \emph{saturada} para esta relación. Ahora bien, como el morfismo $\Spec(A')\to\Spec(A)$ tiene fibras geométricamente conexas (18.9.1), las fibras de $g$ son conexas, de donde se deduce inmediatamente nuestra afirmación.

En particular, si $X$ es localmente conexo (lo que ocurrirá si $X$ es localmente noetheriano), entonces, para toda componente conexa $U$ de $X$ (que es abierta y cerrada

\oldpage[IV-4]{151}

en $X$), $g^{-1}(U)$ es conexa (Bourbaki, \emph{Top. gén.}, cap.~I\textsuperscript{er}, 3\textsuperscript{e} ed., \textsection~11, n\textsuperscript{o}~3, prop.~7). Si se denota por $\pi_0(X)$ el conjunto de las componentes conexas de $X$, la aplicación $\pi_0(X')\to\pi_0(X)$ deducida canónicamente de $g$ es, por tanto, \emph{biyectiva}.

\begin{corollary}[18.9.5]
\label{IV.18.9.5}
Bajo las hipótesis de (18.9.4), el funtor
\begin{equation}
  Z\longmapsto Z\otimes_A A'
  \tag{18.9.5.1}
  \label{IV.18.9.5.1}
\end{equation}
de la categoría de los preesquemas étales sobre $X$ a la de los preesquemas étales sobre $X'$ es plenamente fiel.
\end{corollary}

Sean $Z_1$, $Z_2$ dos preesquemas étales sobre $X$, y pongamos $Z'_i=Z_i\otimes_A A'$ ($i=1,2$); se trata de demostrar que todo $X'$-morfismo $Z'_1\to Z'_2$ procede, por cambio de base, de un único $X$-morfismo $Z_1\to Z_2$. Supongamos primero que $Z_2$ es \emph{separado} sobre $X$; entonces, como $\Hom_X(Z_1,Z_2)$ se identifica con $\Gamma((Z_1\times_X Z_2)/Z_1)$ y $Z_1\times_X Z_2$ es étale y separado sobre $Z_1$, $\Hom_X(Z_1,Z_2)$ se identifica funtorialmente, en virtud de (17.9.3), con el conjunto de las partes abiertas y cerradas $U$ de $Z_1\times_X Z_2$ tales que la restricción a $U$ de la proyección $p:Z_1\times_X Z_2\to Z_1$ sea un morfismo sobreyectivo y radicial. La afirmación se deduce, por tanto, en el caso considerado de (18.9.4) y de que $Z'_1\times_{X'}Z'_2=(Z_1\times_X Z_2)\otimes_A A'$.

Pasemos al caso general: bastará demostrar que el grafo $\Gamma'$ de un $X'$-morfismo $Z'_1\to Z'_2$, que es abierto (17.9.3) en $Z'_1\times_{X'}Z'_2$, es de la forma $\Gamma\otimes_A A'$, donde $\Gamma$ está inducido sobre una parte \emph{abierta} de $Z=Z_1\times_X Z_2$. En efecto, la restricción a $\Gamma$ de la proyección $p:Z\to Z_1$ es entonces un \emph{isomorfismo}, pues, por el cambio de base $A\to A'$, esta restricción se convierte en el isomorfismo dado por la restricción a $\Gamma'$ de la proyección $p':Z'\to Z'_1$ (con $Z'=Z\otimes_A A'$), y basta aplicar (2.7.1, (viii)); la conclusión se deduce entonces de la caracterización de los grafos de morfismos (I, 5.3).

Si $q:Z'\to Z$ es la proyección canónica, para demostrar que $\Gamma'$ es de la forma $q^{-1}(\Gamma)$, donde $\Gamma$ es abierto en $Z$, basta, puesto que $q$ es un morfismo fielmente plano y cuasicompacto, demostrar que existe un conjunto $U\subset Z$ tal que $\Gamma'=q^{-1}(U)$ (2.3.12). Pongamos $Z''=Z'\times_Z Z'=Z\times_X X''$, donde $X''=X'\times_X X'$, y sean $q_1:Z''\to Z'$, $q_2:Z''\to Z'$ las proyecciones canónicas. Aplicando (4.5.19.1), basta mostrar que $q_1^{-1}(\Gamma')=q_2^{-1}(\Gamma')$. Pero esta propiedad es verdadera si y solo si lo es después de cada cambio de base $\Spec(k(x))\to X$, donde $x$ recorre $X$. En otras palabras, basta demostrar el corolario cuando $X$ es el \emph{espectro de un cuerpo}; pero, como todo $X$-preesquema étale es entonces automáticamente \emph{separado} sobre $X$ (17.6.2, $c'$)), se vuelve al caso considerado al comienzo de la demostración.

\begin{remarks}[18.9.6]
\label{IV.18.9.6}
\begin{enumerate}
  \item[(i)] Es posible que, si el cuerpo residual de $A$ es de característica 0, el funtor (18.9.5.1) induzca incluso una \emph{equivalencia} entre la categoría de los recubrimientos étales de $X$ y la categoría de los recubrimientos étales de $X'$. En todo caso, puede demostrarse cuando $A$ es \emph{excelente}, utilizando la resolución de singularidades de Hironaka (M.~Artin). Es plausible que el enunciado análogo siga siendo verdadero sin restricción sobre la característica del cuerpo residual, a condición de limitarse a los recubrimientos «galoisianos principales» cuyo grupo sea de orden primo con la característica residual (cf. [41]).

\oldpage[IV-4]{152}

  \item[(ii)] Se ignora si las fibras formales de $A$ son geométricamente conexas (o incluso geométricamente irreducibles) cuando $A$ es un anillo local henseliano cualquiera. Bastaría con que, para todo anillo local noetheriano, henseliano e \emph{íntegro} $A$, $\Spec(\widehat A)$ fuera irreducible. Se ignora si esto sucede siempre.

  \item[(iii)] Se ignora si, cuando $A$ es un anillo local excelente, su henselizado estricto ${}^{hs}\!A$ es excelente. Puede verse que, para dilucidar esta cuestión, es posible reducirse al caso en que $A=k[[T_1,\ldots,T_n]]$, siendo $k$ un cuerpo de característica $p>0$. Teniendo en cuenta (18.8.17), la cuestión consiste en saber si las fibras formales de ${}^{hs}\!A$ son geométricamente regulares. La respuesta es afirmativa cuando $n=1$, pero no se conoce para $n=2$. Puede demostrarse que la respuesta es afirmativa siempre que, para todo esquema $Y_1$ finito sobre $Y=\Spec(\widehat A)$, puedan resolverse las singularidades de $Y_1$ (7.9.1).
\end{enumerate}
\end{remarks}

La afirmación de conexidad formulada en (18.9.1) se generaliza de la manera siguiente:

\begin{theorem}[18.9.7]
\label{IV.18.9.7}
Sean $A$, $B$ dos anillos locales noetherianos, $\rho:A\to B$ un homomorfismo local, y $f:\Spec(B)=X\to\Spec(A)=Y$ el morfismo correspondiente. Supongamos que se verifican las condiciones siguientes:
\begin{enumerate}
  \item[(i)] $f$ es plano y todas las fibras $X_y=f^{-1}(y)$ ($y\in Y$) son geométricamente reducidas (4.6.2).
  \item[(ii)] O bien el anillo $A$ es geométricamente unirrama (0, 23.2.1), o bien el anillo $A$ es unirrama (0, 23.2.1) y $k(B)$ es una extensión primaria (4.3.1) de $k(A)$.
\end{enumerate}
Entonces, si $\eta$ es el punto genérico de $Y$, $X_\eta$ es conexo y, para toda parte cerrada $T$ de $X$ tal que $f(T)$ sea rara en $Y$, $X-T$ es conexo.
\end{theorem}

Las dos conclusiones enunciadas son, de hecho, equivalentes; en efecto, la hipótesis de que $f(T)$ sea rara implica que $f(T)$ no contiene a $\eta$; por otra parte, como $\{\eta\}$ es proconstructible (I, 9.6) y $f$ es plano, $X_\eta$ es denso en $X$ (2.3.10), y, por consiguiente, si $X_\eta$ es conexo, lo mismo sucede con $X-T$, que lo contiene. Recíprocamente, cuando $U$ recorre el conjunto de las vecindades abiertas afines de $\eta$ en $Y$, $X_\eta$ es límite proyectivo de los subpreesquemas inducidos sobre los abiertos $f^{-1}(U)$ de $X$ (8.1.2, $a)$); si se supone que $X-T$ es conexo para toda parte cerrada $T$ de $X$ tal que $f(T)$ sea rara en $Y$, los $f^{-1}(U)$ son conexos y, por tanto, lo mismo sucede con $X_\eta$ (8.4.1).

Demostremos primero que podemos limitarnos al caso en que $A$ es íntegro y $Y=\Spec(A)$ es geométricamente unirrama (6.15.1) (lo que implica que $A$ es geométricamente unirrama, pero no equivale a esta condición (6.15.2)). Es claro que podemos suponer primero que $Y$ es reducido, considerando el morfismo $X\times_Y Y_{\mathrm{red}}\to Y_{\mathrm{red}}$, deducido de $f$ por cambio de base, que es plano y tiene las mismas fibras que $f$. Podemos suponer, por tanto, que $A$ es íntegro y unirrama; si $K$ es el cuerpo de fracciones de $A$, existe entonces una sub-$A$-álgebra finita $A''$ de $K$ tal que, si $A'$ es la clausura integral de $A$, el morfismo $\Spec(A')\to\Spec(A'')$ sea radicial (0, 23.2.5); se deduce de ello (6.15.5) que $\Spec(A'')$ es geométricamente unirrama. Como $A$ es unirrama, $A'$ es un anillo local y, por tanto, lo mismo sucede con $A''$; demostremos que el anillo $B''=B\otimes_A A''$ también es local. En efecto, $B''$ es una $B$-álgebra finita y, por tanto, un anillo semilocal noetheriano; si $k''$ es el cuerpo residual de $A''$, los ideales maximales de $B''$ son los puntos de $\Spec(B'')$ situados sobre el punto cerrado de $\Spec(B)$ y, por tanto, los puntos de $\Spec(k(B)\otimes_{k(A)}k'')$ (I, 3.4.9)

\oldpage[IV-4]{153}

puesto que el punto cerrado de $\Spec(A'')$ es el único punto situado sobre el punto cerrado de $\Spec(A)$. Ahora bien, cuando $A$ es geométricamente unirrama, $k''$ es una extensión radicial de $k(A)$, por lo que $\Spec(k(B)\otimes_{k(A)}k'')$ es radicial sobre $\Spec(k(B))$ y, por tanto, solo contiene un punto. Cuando $A$ es unirrama y $k(B)$ es una extensión primaria de $k(A)$, $\Spec(k(B)\otimes_{k(A)}k'')$ es irreducible (4.3.2) y, como es un espacio finito y discreto (I, 6.4.4), sigue reduciéndose a un punto, lo que demuestra que, en los dos casos considerados, $B''$ es local. Si $Y''=\Spec(A'')$, $X''=\Spec(B'')=X\times_Y Y''$, el morfismo $f''=f_{(Y'')}:X''\to Y''$ es plano y tiene fibras geométricamente reducidas; además, si $\eta''$ es el punto genérico de $Y''$, $k(\eta'')=k(\eta)=K$ por definición, por lo que las fibras $X_\eta$ y $X''_{\eta''}$ son isomorfas. Basta, por tanto, demostrar que $X''_{\eta''}$ es conexa.

\begin{remark}[18.9.7.1]
\label{IV.18.9.7.1}
Cuando el anillo $A$ es \emph{japonés}, en lo que precede puede tomarse $A''=A'$ por definición (0, 23.1.1); se ve, por tanto, que en este caso incluso se reduciría la demostración del teorema al caso en que $A$ es \emph{íntegramente cerrado}.
\end{remark}

\begin{env}[18.9.7.2]
\label{IV.18.9.7.2}
Supongamos, pues, en adelante que $Y$ es íntegro y geométricamente unirrama. Observemos que, si $T$ es una parte cerrada de $X$ tal que $f(T)$ sea rara en $Y$, $T$ no contiene ningún punto maximal de $X$, puesto que $f$ es plano (2.3.4), y es, por tanto, rara en $X$; como $X$ es un esquema local y, por consiguiente, conexo, para demostrar que $X-T$ es conexo basta, en virtud de (15.5.6.1), mostrar que, para todo $x\in X$ tal que $f(x)\neq\eta$, $\Spec(\sh{O}_{X,x})-\{x\}$ es conexo. Pongamos $y=f(x)$, $A_1=\sh{O}_{Y,y}$, $B_1=\sh{O}_{X,x}$, $Y_1=\Spec(A_1)$, $X_1=\Spec(B_1)$, y sea $f_1:X_1\to Y_1$ el morfismo correspondiente al homomorfismo local $A_1\to B_1$ deducido de $\rho$; es claro que $f_1$ es plano, y se deduce de (4.6.1) que sus fibras son geométricamente reducidas; además, $A_1$ es íntegro y geométricamente unirrama (6.15.1), y $\dim(A_1)\geq 1$. Se reduce así la cuestión a demostrar el
\end{env}

\begin{lemma}[18.9.7.3]
\label{IV.18.9.7.3}
Sean $A$, $B$ dos anillos locales noetherianos, $Y=\Spec(A)$, $X=\Spec(B)$, $\rho:A\to B$ un homomorfismo local, y $f:X\to Y$ el morfismo correspondiente. Se supone que se verifican las condiciones siguientes:
\begin{enumerate}
  \item[(i)] $f$ es plano y todas las fibras $X_y$ ($y\in Y$) son geométricamente reducidas.
  \item[(ii)] $A$ es íntegro, geométricamente unirrama y $\dim(A)\geq 1$.
\end{enumerate}
Entonces, si $b$ es el punto cerrado de $X$, $X-\{b\}$ es conexo.
\end{lemma}

Observemos primero que, en virtud del teorema de Hartshorne (5.10.7), $X-\{b\}$ es conexo si se tiene $\operatorname{prof}(B)\geq 2$. Ahora bien, se tiene (6.3.1)
\begin{equation}
  \operatorname{prof}(B)=\operatorname{prof}(A)+\operatorname{prof}(B\otimes_A k)
  \tag{18.9.7.4}
  \label{IV.18.9.7.4}
\end{equation}
denotando por $k$ el cuerpo residual de $A$. Por otra parte, puesto que $A$ es íntegro y $\dim(A)\geq 1$, se tiene $\operatorname{prof}(A)\geq 1$ y, por tanto, $\operatorname{prof}(B)\geq 2$ \emph{salvo} cuando $\operatorname{prof}(B\otimes_A k)=0$; además, $B\otimes_A k$ es reducido por (i), de modo que la relación $\operatorname{prof}(B\otimes_A k)=0$ significa que $B\otimes_A k$ es un \emph{cuerpo} (0, 16.4.7). Supondremos, pues, en adelante que se verifica esta última condición. Comenzaremos tratando un caso en el que la demostración es muy sencilla.

A) \emph{Caso en que $A$ es íntegramente cerrado.} Entonces, si $\dim(A)\geq 2$, se tiene $\operatorname{prof}(A)\geq 2$ (0, 16.5.1) y, por tanto, $\operatorname{prof}(B)\geq 2$. Solo queda, pues, considerar el caso en que $\dim(A)=1$ y $B\otimes_A k$ es un cuerpo; $A$ es entonces un anillo de valoración discreta, y por tanto regular, y puesto que $B\otimes_A k$

\oldpage[IV-4]{154}

es un cuerpo y $f$ es plano, $B$ es regular (0, 17.3.3); pero, además, se tiene (6.1.1.1) $\dim(B)=\dim(A)+\dim(B\otimes_A k)=\dim(A)=1$, por lo que $B$ es también un anillo de valoración discreta. Pero entonces $X-\{b\}$ se reduce a un solo punto, de donde se deduce el lema (18.9.7.3) en este caso.

Se observará que esto demuestra el enunciado del teorema (18.9.7) cuando se supone además que, en dicho enunciado, el anillo $A$ es \emph{japonés}, pues, en virtud de la observación (18.9.7.1), se vuelve al caso en que $A$ es \emph{íntegramente cerrado} y entonces, en la reducción precedente (18.9.7.3), el anillo $A_1=\sh{O}_{Y,y}$ es también íntegramente cerrado y puede aplicarse el resultado que acaba de demostrarse.

B) \emph{Caso en que $\dim(A)\geq 2$.} El resultado de (18.9.7.3) se deducirá entonces de los dos lemas siguientes:

\begin{lemma}[18.9.7.5]
\label{IV.18.9.7.5}
Sean $A$ un anillo semilocal noetheriano y reducido, $A'$ su clausura integral en su anillo total de fracciones, $X=\Spec(A)$, $X'=\Spec(A')$; sean $S$ el conjunto de los puntos cerrados de $X$ y $U=X-S$. Supongamos que, para todo punto $x'\in X'$ situado sobre un punto de $S$, se tiene $\dim(\sh{O}_{X',x'})\geq 2$; entonces el anillo $A^{(1)}=\Gamma(U,\sh{O}_X)$ es entero sobre $A$ y es una $A$-álgebra isomorfa a una subálgebra de $A'$.
\end{lemma}

Recordemos que, si los $\mathfrak p_i$ ($1\leq i\leq m$) son los ideales primos minimales de $A$, $A'$ es el producto directo de las clausuras integrales $A'_i$ de los anillos íntegros $A_i=A/\mathfrak p_i$, de modo que $X'$ es la suma de los esquemas $X'_i=\Spec(A'_i)$. Si $X_0$ es la suma de los esquemas $X_i=\Spec(A_i)$ ($1\leq i\leq m$) y $U_0$ la imagen inversa de $U$ en $X_0$, el homomorfismo canónico $\Gamma(U,\sh{O}_X)\to\Gamma(U_0,\sh{O}_{X_0})$ es inyectivo, puesto que $X$ es reducido y el morfismo $X_0\to X$ (y, por tanto, también el morfismo $U_0\to U$) es sobreyectivo. Como $\Gamma(U_0,\sh{O}_{X_0})$ es el producto directo de los $\Gamma(U_i,\sh{O}_{X_i})$ ($1\leq i\leq m$), donde $U_i$ es la imagen inversa de $U$ en $X_i$, todo se reduce a demostrar que, para cada $i$, $\Gamma(U_i,\sh{O}_{X_i})$ es isomorfo a una sub-$A_i$-álgebra de $A'_i$. En otras palabras, se vuelve al caso en que $A$ es un anillo semilocal noetheriano íntegro. Como $X$ es reducido y el morfismo $X'\to X$ es sobreyectivo, el homomorfismo $\Gamma(U,\sh{O}_X)\to\Gamma(U',\sh{O}_{X'})$ (donde $U'$ es la imagen inversa de $U$ en $X'$) es inyectivo, y todo se reduce a ver que el homomorfismo canónico $A'=\Gamma(X',\sh{O}_{X'})\to\Gamma(U',\sh{O}_{X'})$ es \emph{biyectivo}. Ahora bien, por (I, 8.2.1.1), $\Gamma(U',\sh{O}_{X'})$ es la intersección de los anillos locales $A'_{\mathfrak p'}$, donde $\mathfrak p'$ recorre $U'$, y, en virtud de la hipótesis, entre estos anillos locales figuran todos aquellos para los que $\mathfrak p'$ es de altura 1. Pero el razonamiento de (0, 23.2.7) se aplica también a un anillo semilocal íntegro noetheriano, por lo que $A'$ es un \emph{anillo de Krull semilocal} y es, por tanto, la intersección de los anillos locales $A'_{\mathfrak p'}$, donde $\mathfrak p'$ recorre el conjunto de los ideales primos de altura 1 de $A'$ (Bourbaki, \emph{Alg. comm.}, cap.~VII, \textsection~1, n\textsuperscript{o}~6, th.~4); \emph{a fortiori}, $A'$ es la intersección de los $A'_{\mathfrak p'}$ para $\mathfrak p'\in U'$, lo que termina la demostración del lema.

\begin{remark}[18.9.7.6]
\label{IV.18.9.7.6}
Se sabe por (0, 23.2.5) que existe una sub-$A$-álgebra finita $A''$ del cuerpo de fracciones $K$ de $A$ tal que el morfismo $\Spec(A')\to\Spec(A'')$ sea \emph{radicial}; como este morfismo es también entero y dominante, y por tanto sobreyectivo (II, 6.1.10), es un \emph{homeomorfismo} (2.4.5). En virtud de la interpretación geométrica de $\dim(\sh{O}_{X',x'})$ (5.1.2), se tiene, por tanto, para todo punto $x''$ de $X''=\Spec(A'')$ situado sobre un punto $x\in X$, $\dim(\sh{O}_{X'',x''})=\dim(\sh{O}_{X',x'})\geq 2$ si $x'$ es el único punto de $X'$ situado sobre $x$. Dicho

\oldpage[IV-4]{155}

esto, se deduce en primer lugar de (0, 16.1.5) y (0, 16.1.4.1) que la relación $\dim(\sh{O}_{X',x'})\geq 2$ implica $\dim(\sh{O}_{X,x})\geq 2$. Recíprocamente, supongamos que $\dim(\sh{O}_{X,x})\geq 2$ y que $A$ es universalmente catenario (5.6.2); entonces lo mismo sucede con $\sh{O}_{X,x}$ (5.6.3), y por tanto $\dim(\sh{O}_{X'',x''})=\dim(\sh{O}_{X,x})\geq 2$ en virtud de (5.6.10), de modo que se tiene $\dim(\sh{O}_{X',x'})\geq 2$. La misma conclusión es válida si $\sh{O}_{X,x}$ es \emph{geométricamente unirrama}, pues la fibra del morfismo $X'\to X$ en el punto $x$ se reduce entonces a un solo punto, y $\sh{O}_{X',x'}$ es una $\sh{O}_{X,x}$-álgebra \emph{entera}; basta, pues, aplicar (0, 16.1.5).
\end{remark}

\begin{lemma}[18.9.7.7]
\label{IV.18.9.7.7}
Sean $A$ un anillo local noetheriano, de dimensión $\geq 2$, íntegro y geométricamente unirrama, $Y=\Spec(A)$, e $y$ el punto cerrado de $Y$. Sean $X$ un preesquema localmente noetheriano y conexo, y $f:X\to Y$ un morfismo plano. Entonces, para toda parte cerrada $T\subset f^{-1}(y)$, $X-T$ es conexo.
\end{lemma}

El mismo razonamiento que al comienzo de (18.9.7.2) reduce la cuestión a demostrar que, para todo punto $x\in f^{-1}(y)$, $\Spec(\sh{O}_{X,x})-\{x\}$ es conexo. Podemos, por tanto, limitarnos al caso en que $X=\Spec(B)$, donde $B$ es un anillo local noetheriano y el homomorfismo $A\to B$ correspondiente a $f$ es local; pongamos $U=Y-\{y\}$ y observemos que $V=f^{-1}(U)$ es denso en $X$ (2.3.10); basta, por tanto, demostrar que $V$ es \emph{conexo}. Conservemos las notaciones de (18.9.7.5) y pongamos, por tanto, $A^{(1)}=\Gamma(U,\sh{O}_Y)$, $Y_1=\Spec(A^{(1)})$, $X_1=X\times_Y Y_1=\Spec(B\otimes_A A^{(1)})$, de modo que se tiene el diagrama conmutativo
\[
  \xymatrix{
    X \ar[d]_{f} & X_1 \ar[l]^{g_1} \ar[d]^{f_1} \\
    Y & Y_1 \ar[l]_{g}
  }
\]

En virtud de (2.3.1) y de la definición de $A^{(1)}$, se tiene un isomorfismo canónico $\Gamma(V,\sh{O}_X)\xrightarrow{\sim}B\otimes_A A^{(1)}$, y basta, por tanto, demostrar que este último anillo es \emph{local}, pues tal anillo no puede ser producto de dos anillos no reducidos a 0 (III, 7.8.6.1). Pero se ha visto en (18.9.7.5) que $A^{(1)}$ es una sub-$A$-álgebra de la clausura integral $A'$ de $A$, puesto que $\dim(A)\geq 2$. La hipótesis de que $A$ es geométricamente unirrama implica que el morfismo $\Spec(A')\to\Spec(A)$ es radicial en el punto $y$ (6.15.3); \emph{a fortiori}, $g:Y_1\to Y$ es radicial en el punto $y$ y, por consiguiente, $g_1$ es radicial en el punto cerrado $x$ de $X$; como, además, $g_1$ es un morfismo entero, los puntos cerrados de $X_1$ están situados sobre $x$, por lo que solo existe un punto cerrado de $X_1$, lo que termina la demostración de (18.9.7.7).

Es claro que (18.9.7.7) demuestra (18.9.7.3) cuando $\dim(A)\geq 2$ (incluso sin suponer que las fibras $X_y$ sean geométricamente reducidas).

C) \emph{Caso en que $\dim(A)=1$.} Se ha visto al comienzo de la demostración de (18.9.7.3) que puede suponerse $\dim(B\otimes_A k)=0$ y, por planitud (6.1.1.1), se tiene, por tanto,
\[
  \dim(B)=\dim(A)+\dim(B\otimes_A k)=1.
\]
Como $A$ es íntegro, $Y$ se compone de dos puntos, el punto genérico $\eta$ y el punto cerrado $a$. Además, puesto que $f$ es plano, toda componente irreducible de $X$ domina $Y$ (2.3.4), por lo que $X-\{b\}$, que es el conjunto de los puntos maximales $\xi_i$ de $X$, es igual al

\oldpage[IV-4]{156}

conjunto subyacente a $f^{-1}(\eta)$, y la fibra $f^{-1}(\eta)$ es suma de los $\Spec(k(\xi_i))$. La hipótesis sobre los $X_y$ y el hecho de que estas fibras sean de dimensión 0 muestran que son \emph{geométricamente regulares} (6.7.6) y, \emph{a fortiori}, geométricamente normales; dicho de otro modo, $f$ es un \emph{morfismo normal}. Pero, como $A$ es íntegro y geométricamente unirrama, se deduce entonces de (6.15.10) que $B$ es íntegro y geométricamente unirrama, por lo que $f^{-1}(\eta)=X-\{b\}$ se reduce a un solo punto y, \emph{a fortiori}, es conexo; esto termina, por tanto, la demostración de (18.9.7.3) y la de (18.9.7).

\begin{remark}[18.9.7.8]
\label{IV.18.9.7.8}
En el caso C) de la demostración de (18.9.7.3), puede evitarse recurrir al delicado resultado (6.15.10) razonando de la manera siguiente: como $A$ es íntegro y de dimensión 1, se deduce del teorema de Krull--Akizuki (Bourbaki, \emph{Alg. comm.}, cap.~VII, \textsection~2, n\textsuperscript{o}~5, prop.~5) que su clausura integral $A'$ es un anillo \emph{noetheriano}. El mismo razonamiento que al comienzo de la demostración de (18.9.7) muestra entonces que $B'=B\otimes_A A'$ es un anillo local; además, por planitud, $B'$ está contenido en el anillo total de fracciones $R$ del anillo reducido $B$ (3.3.5). Ahora bien, el razonamiento que demuestra el teorema de Krull--Akizuki (Bourbaki, \emph{loc. cit.}) se aplica igualmente a un anillo local noetheriano \emph{reducido} de dimensión 1 y muestra que, para tal anillo $B$, todo anillo comprendido entre $B$ y su anillo total de fracciones es \emph{noetheriano}. Como el anillo $B'$ es un anillo local noetheriano y el morfismo $\Spec(B')\to\Spec(A')$ es normal, $B'$ es un anillo íntegro e íntegramente cerrado (6.5.4), y se concluye como al comienzo de la demostración de (18.9.7).
\end{remark}

\begin{corollary}[18.9.8]
\label{IV.18.9.8}
Con las notaciones de (18.9.7), supongamos que se verifica la condición (i) de (18.9.7) y una de las dos condiciones siguientes:
\begin{enumerate}
  \item[(ii$'$)] $A$ es un anillo estrictamente local (18.8.2).
  \item[(ii$''$)] $A$ es un anillo local henseliano y $k(B)$ es una extensión primaria de $k(A)$.
\end{enumerate}
Entonces todas las fibras $X_y$ ($y\in Y$) son geométricamente conexas.
\end{corollary}

Sea $\mathfrak p$ el ideal primo de $A$ correspondiente al punto $y$, y pongamos $A'=A/\mathfrak p$ y $B'=B\otimes_A A'=B/\mathfrak pB$, que son evidentemente anillos locales noetherianos; es claro que el morfismo $f':\Spec(B')\to\Spec(A')$ verifica la condición (i) de (18.9.7); como $y$ es el punto genérico de $\Spec(A')$ y $k(A')=k(A)$, $k(B')=k(B)$, se ve que basta verificar que, en el caso (ii$'$) (resp. (ii$''$)), $A'$ es geométricamente unirrama (resp. unirrama). Pero $A'$ es íntegro y, en el caso (ii$'$) (resp. (ii$''$)), es estrictamente local (resp. henseliano) en virtud de (18.5.10). Se concluye que $A'$ es geométricamente unirrama (resp. unirrama) en virtud de (18.8.15) (resp. (18.6.12)).

\begin{corollary}[18.9.9]
\label{IV.18.9.9}
Sea $A$ un anillo local noetheriano henseliano, universalmente japonés (\emph{i.e.} (7.6.4) y (7.7.2)), cuyas fibras formales (7.3.13) son geométricamente reducidas; entonces las fibras formales de $A$ son geométricamente conexas.
\end{corollary}

Basta aplicar (18.9.8) al caso en que $B=\widehat A$, puesto que $B$ es un $A$-módulo plano y que $k(B)=k(A)$.

\begin{remark}[18.9.10]
\label{IV.18.9.10}
Las demostraciones de (18.9.4) y (18.9.5) solo utilizan el hecho de que las fibras formales de $A$ son geométricamente conexas. Las conclusiones

\oldpage[IV-4]{157}

de estas dos proposiciones siguen siendo, por tanto, válidas cuando se supone que el anillo local $A$ es noetheriano, henseliano y universalmente japonés.
\end{remark}

\begin{corollary}[18.9.11]
\label{IV.18.9.11}
Sean $X$, $Y$ dos preesquemas localmente noetherianos, suponiéndose que $Y$ es geométricamente unirrama y que $X$ es conexo. Sea $f:X\to Y$ un morfismo reducido (\emph{i.e.}, plano y de fibras geométricamente reducidas (6.8.1)). Entonces, para toda parte cerrada $T$ de $X$ tal que $f(T)$ sea rara en $Y$, $X-T$ es conexo.
\end{corollary}

Como $f$ es plano y $f(T)$, al ser rara, no puede contener ningún punto maximal de $Y$, $T$ no puede contener ningún punto maximal de $X$ (2.3.4) y es, por tanto, rara en $X$. Utilizando (15.5.6.1), basta demostrar que, para todo $x\in T$, $\Spec(\sh{O}_{X,x})-\{x\}$ es conexo. Pongamos $y=f(x)$, que no es un punto maximal de $Y$. Como el anillo $\sh{O}_{Y,y}$ es geométricamente unirrama, puede aplicarse (18.9.7) al morfismo
\[
  f_1:\Spec(\sh{O}_{X,x})\longrightarrow\Spec(\sh{O}_{Y,y}),
\]
que verifica las condiciones (i) y (ii) de este teorema, y a la parte cerrada $\{x\}$ de $\Spec(\sh{O}_{X,x})$,
puesto que $f_1(\{x\})=\{y\}$ es rara en $\Spec(\sh{O}_{Y,y})$.
